\chapter{Interrupt Timing and Scheduler Response}
\section{Stress Response-Latency Distribution}
An interrupt that is taken on time has still done nothing useful. The work begins when the
dispatcher task runs, and the interval between the two is where the scheduler makes itself felt.
For event $i$ the firmware records the cycle counter at handler entry and again when the
dispatcher removes that event from the ring, so that
\begin{equation}
  L_i \;=\; \frac{c^{\mathrm{disp}}_i - c^{\mathrm{isr}}_i}{f_{\mathrm{cpu}}},
  \qquad f_{\mathrm{cpu}} = \SI{160}{\mega\hertz}.
\end{equation}
With the worker parked, $L_i$ sits at \SI{13.2}{\micro\second} and barely moves. With an
equal-priority worker on the same core, the mean rises to \SI{53.7}{\micro\second} and more than
half the events wait longer than the period that generated them (\cref{fig:ecdf}).

\begin{figure}[htbp]
\centering
\begin{tikzpicture}
\begin{axis}[rp axis, width=0.78\linewidth, height=5.4cm,
  xlabel={Dispatch latency $L$ (\si{\micro\second})}, ylabel={Fraction of events $\le L$},
  xmin=0, xmax=125, ymin=0, ymax=1.02, legend pos=south east]
\addplot[rp base] table[x=x,y=p]{../build/data/lat_ecdf_baseline.dat}; \addlegendentry{baseline}
\addplot[rp accent] table[x=x,y=p]{../build/data/lat_ecdf_stress.dat}; \addlegendentry{worker at priority 22}
\draw[rp ref] (axis cs:50,0) -- (axis cs:50,1.02) node[pos=0.93, anchor=west, rp note]{\SI{50}{\micro\second} period};
\end{axis}
\end{tikzpicture}
\caption{\evM\ Empirical distribution of dispatch latency, original firmware, ten runs per
condition (\num{13417} and \num{13715} events). The latency is handler entry to dispatcher pop, not a
context-switch time; crossing the \SI{50}{\micro\second} line is a period exceedance, not a missed deadline.}
\label{fig:ecdf}
\end{figure}

\begin{table}[htbp]
\centering
\caption{\evM\ Dispatch latency by condition. Run-level range in brackets; event-level SD is descriptive.}
\label{tab:lat}
\begin{tabular}{l S[table-format=2.3] S[table-format=2.3] S[table-format=3.3] S[table-format=2.1] l}
\toprule
Condition & {Mean} & {SD} & {P99} & {$>\SI{50}{\micro\second}$} & Run means \\
          & {(\si{\micro\second})} & {(\si{\micro\second})} & {(\si{\micro\second})} & {(\si{\percent})} & (\si{\micro\second}) \\
\midrule
Baseline            & 13.224 & 1.268 & 21.094 & 0.0  & [13.22,\,13.24] \\
Worker, priority 22 & 53.664 & 33.237 & 113.169 & 53.3 & [53.15,\,54.25] \\
\bottomrule
\end{tabular}
\end{table}

The mechanism is in six lines of the worker: at priority 22 it yields only after a full chunk of
work, and until it does the dispatcher is ready but not running.
\begin{lstlisting}[style=fw, firstnumber=856, title={\lsttitle{esp32\_benchmark\_logged.ino}{856--861}}]
        while (bench_stress_enabled) {
            bench_run_stress_chunk(state);
            if (bench_stress_priority == 22U) {
                taskYIELD();
            } else if (bench_stress_priority > 22U) {
                const uint32_t elapsed = bench_read_ccount() - burst_start;
\end{lstlisting}
